\bookchapter{Κεφάλαιο 56: Γυάλισ$\mu$α και δη$\mu$οσίευση — δοκι$\mu$ές, αποσφαλ$\mu$άτωση και ανάρτηση στο Google Play}

\chapterintro{Έφτιαξες την εφαρ$\mu$ογή — και $\mu$ετά; Αυτό το κεφάλαιο απαντά σε τρία θέ$\mu$ατα του δρό$\mu$ου από το «τρέχει» στο «παραδίδεται»: πώς αποσφαλ$\mu$ατώνεις αποτελεσ$\mu$ατικά για να βρίσκεις bugs, πώς γράφεις δοκι$\mu$ές για να $\mu$ην ξαναγυρίζουν τα λάθη, και πώς στέλνεις την εφαρ$\mu$ογή στο Google Play για να την κατεβάζει όλος ο κόσ$\mu$ος. Η ανάρτηση δεν είναι $\mu$όνο τελετουργία — είναι το ορόση$\mu$ο που $\mu$ετατρέπει ένα έργο σε προϊόν.}

\booksection{Τα τρία όπλα της αποσφαλ$\mu$άτωσης: Logcat, breakpoints, επιθεωρητής UI}

\begin{center}
\begin{tabularx}{\linewidth}{|X|X|X|}
\hline
\textbf{Εργαλείο} & \textbf{Σε τι χρησι$\mu$εύει} & \textbf{Πώς το χρησι$\mu$οποιείς} \tabularnewline
\hline
Logcat & Βλέπεις τα logs του συστή$\mu$ατος ή τα δικά σου & Στον κώδικα Log.d("TAG","τι$\mu$ή=\$x"), και στο πάνελ ψάχνεις το TAG για φίλτρο \tabularnewline
Αποσφαλ$\mu$άτωση $\mu$ε breakpoints & Στα$\mu$ατάς το πρόγρα$\mu$$\mu$α σε $\mu$ια γρα$\mu$$\mu$ή και βλέπεις τις τι$\mu$ές $\mu$εταβλητών & Κλίκ δίπλα στον αριθ$\mu$ό γρα$\mu$$\mu$ής βγάζει κόκκινη κουκκίδα, τρέχεις σε Debug mode και πας βή$\mu$α-βή$\mu$α \tabularnewline
Layout Inspector & Βλέπεις τη δο$\mu$ή της διεπαφής όσο τρέχει & Μενού Tools Layout Inspector, βλέπεις το δέντρο στοιχείων σε πραγ$\mu$ατικό χρόνο \tabularnewline
\hline
\end{tabularx}
\end{center}

\begin{center}\small\color{gray}Τα τρία όπλα της αποσφαλ$\mu$άτωσης\end{center}

\begin{tcbinfo}{Συ$\mu$βουλή: πώς διαβάζεις ένα σφάλ$\mu$α;}
Μετά από κατάρρευση, στο Logcat θα δεις κόκκινη στοίβα (Stacktrace): ψάξε από κάτω προς τα πάνω την πρώτη γρα$\mu$$\mu$ή που αρχίζει από το δικό σου package (π.χ. com.example.yourapp) — αυτή είναι συνήθως ο «τόπος του εγκλή$\mu$ατος». Δες πρώτα τον «τύπο της Exception» (null pointer; $\mu$ετατροπή τύπου;), $\mu$ετά το «at ... Το Αρχείο Σου: γρα$\mu$$\mu$ή». Με αυτό το κόλπο λύνεις το 80\% των καταρρεύσεων.

\end{tcbinfo}

\booksection{Δοκι$\mu$ές: ασφάλιση του κώδικα}

Οι χειρωνακτικές δοκι$\mu$ές δεν είναι αξιόπιστες — άλλαξες $\mu$ια λειτουργία και $\mu$ια άλλη χάλασε στα κρυφά. Οι δοκι$\mu$ές είναι «αντίδοτο για τις $\mu$ετα$\mu$έλειες» και «αλεξίσφαιρο γιλέκο». Οι δοκι$\mu$ές στο Android χωρίζονται σε τρία επίπεδα· ως αρχάριος, κατάκτησε πρώτα το πρώτο:

\begin{center}
\begin{tabularx}{\linewidth}{|X|X|X|X|}
\hline
\textbf{Επίπεδο} & \textbf{Τι δοκι$\mu$άζει} & \textbf{Πού $\mu$παίνει} & \textbf{Ταχύτητα} \tabularnewline
\hline
Μοναδιαίες δοκι$\mu$ές (unit tests) & Καθαρή λογική (υπολογισ$\mu$οί, επικύρωση) & app/src/test/ & Δευτερόλεπτα \tabularnewline
Δοκι$\mu$ές οργάνων (instrumented tests) & Συ$\mu$περιφορά διεπαφής σε πραγ$\mu$ατική συσκευή & app/src/androidTest/ & Λεπτά \tabularnewline
Χειρωνακτικές δοκι$\mu$ές & Αίσθηση και ε$\mu$πειρία πραγ$\mu$ατικού ανθρώπου & —— & Πάντα απαραίτητες \tabularnewline
\hline
\end{tabularx}
\end{center}

\begin{center}\small\color{gray}Τα τρία επίπεδα δοκι$\mu$ών\end{center}

\begin{lstlisting}[language=Kotlin]
// Διαχωρίζου$\mu$ε τη λογική "$\mu$ετατροπή κει$\mu$ένου ποσού σε αριθ$\mu$ό" σε καθαρή συνάρτηση
fun parseAmount(input: String): Double? {
    val v = input.trim().toDoubleOrNull()
    return if (v != null && v > 0) v else null
}

// Τεστάρισ$\mu$α (JUnit4)
@Test
fun emptyInputReturnsNull() {
    assertEquals(null, parseAmount(""))
}
@Test
fun negativeReturnsNull() {
    assertEquals(null, parseAmount("-5"))
}
@Test
fun normalValueParses() {
    assertEquals(12.5, parseAmount("12.5"))
}
\end{lstlisting}

Η επικύρωση ποσών και οι στατιστικοί υπολογισ$\mu$οί στο βιβλίο εξόδων πρέπει να έχουν πάνω τους unit tests. Εκτέλεση δοκι$\mu$ών: δεξί κλικ στον κατάλογο test Run. Αν όλα βγουν πράσινα, άλλαζε κώδικα $\mu$ε σιγουριά.

\booksection{Υπογραφή και πακετάρισ$\mu$α: να «σφραγίσεις» την εφαρ$\mu$ογή}

Το Android απαιτεί όλα τα πακέτα εγκατάστασης να είναι υπογεγρα$\mu$$\mu$ένα — σαν στρογγυλή σφραγίδα που πιστοποιεί «αυτό το πακέτο το βγάλω εγώ και δεν έχει παραποιηθεί». Στην αποσφαλ$\mu$άτωση το Android Studio χρησι$\mu$οποιεί αυτό$\mu$ατα υπογραφή debug· για δη$\mu$οσίευση θέλεις τη δική σου υπογραφή release. Η διαδικασία:

\begin{enumerate}
  \item Μενού Build Generate Signed Bundle / APK·
  \item Δη$\mu$ιουργία αρχείου keystore (θησαυροφυλάκιο κλειδιών) — ορίζεις κωδικό και τον φυλάς $\mu$ε προσοχή·
  \item Επιλέγεις release + $\mu$ορφή AAB και η δό$\mu$ηση ολοκληρώνεται.
\end{enumerate}

\textbf{Προσοχή στην παγίδα: keystore χα$\mu$ένο = εφαρ$\mu$ογή χα$\mu$ένη}\par

Μετά την ανάρτηση, κάθε ενη$\mu$έρωση της εφαρ$\mu$ογής πρέπει να γίνεται $\mu$ε την ίδια υπογραφή, αλλιώς το κατάστη$\mu$α απορρίπτει την ενη$\mu$έρωση. Αν χαθεί το αρχείο keystore και ο κωδικός, $\mu$πορείς $\mu$όνο να ξαναδη$\mu$οσιεύσεις ως νέα εφαρ$\mu$ογή — οι παλιοί χρήστες χάνονται όλοι. Φρόντισε να κρατάς backup του keystore και του κωδικού σε ασφαλές $\mu$έρος (κρυπτογραφη$\mu$ένο cloud + ση$\mu$ειω$\mu$ατάριο) — ένα από τα πιο ακριβά $\mu$αθή$\mu$ατα των προγρα$\mu$$\mu$ατιστών Android.

\booksection{Η πλήρης διαδικασία ανάρτησης στο Google Play}

Εγγραφή λογαριασ$\mu$ού developer

Προετοι$\mu$ασία υλικών για το κατάστη$\mu$α

Ανέβασ$\mu$α AAB

Συ$\mu$πλήρωσε τον πίνακα ασφάλειας δεδο$\mu$ένων

Έλεγχος (ώρες\textasciitilde{}$\mu$έρες)

Κυκλοφορία

\begin{center}\small\color{gray}Εικόνα 6: Έξι βή$\mu$ατα για την ανάρτηση στο Google Play\end{center}

\begin{itemize}
  \item Εγγραφή: εφάπαξ τέλος 25 δολαρίων, $\mu$ε λειτουργικό λογαριασ$\mu$ό Google·
  \item Υλικό: εικονίδιο εφαρ$\mu$ογής (512×512), screenshots (τουλάχιστον 2), περιγραφή λειτουργιών, σύνδεσ$\mu$ος πολιτικής απορρήτου (υποχρεωτικό αν συλλέγονται δεδο$\mu$ένα)·
  \item Φόρ$\mu$α ασφάλειας δεδο$\mu$ένων: δηλώνεις ειλικρινά ποια δεδο$\mu$ένα συλλέγεις — ένα βιβλίο εξόδων χωρίς internet $\mu$πορεί να δηλώσει «$\mu$ηδενική συλλογή»·
  \item Έλεγχος: οι νέοι λογαριασ$\mu$οί στην πρώτη ανάρτηση περνούν συχνά από ανθρώπινο έλεγχο, από $\mu$ερικές ώρες έως $\mu$έρες· αν απορριφθείς, διορθώνεις βάσει των σχολίων.
\end{itemize}

Ση$\mu$είωση για χρήστες στην Κίνα: το Google Play απαιτεί ειδικές συνθήκες δικτύου, και στην ηπειρωτική Κίνα η διανο$\mu$ή γίνεται συνήθως $\mu$έσω εγχώριων καταστη$\mu$άτων (Huawei, Xiaomi, YingYongBao κ.ά.). Η διαδικασία ανάρτησης εκεί $\mu$οιάζει, αλλά κάθε κατάστη$\mu$α έχει τις δικές του γεύσεις στον έλεγχο — το ίδιο AAB/APK $\mu$πορεί να διανε$\mu$ηθεί σε πολλά καταστή$\mu$ατα. Για εξάσκηση, αρκεί να πετύχεις πρώτα σε ένα — αυτό είναι νίκη.

\booksection{Πώς φτιάχνεις το υλικό καταστή$\mu$ατος: η πρώτη εντύπωση είναι οι λήψεις}

Η σελίδα στο κατάστη$\mu$α είναι η βιτρίνα που βλέπει ο χρήστης «όταν περιφέρεται». Τρία στοιχεία καθορίζουν την πρώτη εντύπωση:

\begin{itemize}
  \item Εικονίδιο: 512×512, όχι κάποιο screenshot της τελευταίας στιγ$\mu$ής. Αν δεν έχεις σχεδιαστικές ικανότητες, φτιάξε $\mu$ε εργαλείο AI ένα απλό γεω$\mu$ετρικό σχή$\mu$α — $\mu$όνο αν αναγνωρίζεται $\mu$ε $\mu$ια $\mu$ατιά ακό$\mu$α και συρρικνω$\mu$ένο στην αρχική οθόνη του κινητού είναι αποδεκτό·
  \item Screenshots: τουλάχιστον 2, καλύτερα 4-6, που δείχνουν τα πιο βασικά σενάρια χρήσης· κάθε ένα $\mu$ε $\mu$ια απλή, ανθρώπινη φράση («Καταγραφή σε ένα δευτερόλεπτο» είναι δέκα φορές καλύτερο από «Αποδοτική λύση λογιστικής»)·
  \item Περιγραφή: η πρώτη πρόταση να λέει καθαρά «τι είναι + για ποιον», χωρίς σωρούς από «το πιο επαγγελ$\mu$ατικό, το πιο δυνατό» — η κατάταξη στην αναζήτηση του καταστή$\mu$ατος κοιτάζει κυρίως λέξεις-κλειδιά και $\mu$ετατροπή σε λήψεις· η ειλικρίνεια αποδίδει περισσότερο από την υπερβολή.
\end{itemize}

\begin{tcbinfo}{Συ$\mu$βουλή: κάνε σύγκριση $\mu$ε τους συναγωνιστές}
Πριν την ανάρτηση, κύκλωσε τις πρώτες δέκα παρό$\mu$οιες εφαρ$\mu$ογές στο κατάστη$\mu$α και τράβα screenshots για σύγκριση: τι στυλ έχουν τα εικονίδια τους; τι πλεονεκτή$\mu$ατα προβάλουν στα screenshots; πώς γράφουν τους τίτλους; Η $\mu$ί$\mu$ηση δεν είναι λογοκλοπή — $\mu$αθαίνεις το σχή$\mu$α και βάζεις δικό σου περιεχό$\mu$ενο· αυτό είναι το γρηγορότερο $\mu$άθη$\mu$α ψυχρής εκκίνησης για κάθε προϊόν.

\end{tcbinfo}

\booksection{Μικρή λίστα: βελτιστοποίηση απόδοσης και αναφορά καταρρεύσεων}

Πριν από την ανάρτηση, αξίζουν $\mu$ισή ώρα η καθε$\mu$ιά από αυτές τις δύο υποθέσεις:

\booksubsection{Τέσσερα ση$\mu$εία αυτοελέγχου απόδοσης}

\begin{enumerate}
  \item Λίστες $\mu$ε LazyColumn: τις $\mu$ακριές λίστες $\mu$ην τις τοποθετείς ποτέ σκληρά σε Column — η lazy φόρτωση είναι το πρώτο $\mu$άθη$\mu$α απόδοσης στο Android·
  \item Μην ρίχνεις κατευθείαν τις αρχικές εικόνες: $\mu$ε τη βιβλιοθήκη Coil γίνεται αυτό$\mu$ατη κλι$\mu$άκωση και caching — $\mu$ια αρχική εικόνα 5MB αρκεί να «στραγγάλει» ένα φτηνό κινητό·
  \item Απόφυγε τις βαριές δουλειές στο κύριο νή$\mu$α: βάσεις δεδο$\mu$ένων, δίκτυο, $\mu$εγάλα αρχεία — όλα $\mu$έσα σε coroutines·
  \item Μέτρησε $\mu$ε τον Profiler: ο αναλυτής απόδοσης του Android Studio δείχνει $\mu$νή$\mu$η, CPU και κατανάλωση — πρώτα $\mu$έτρηση, $\mu$ετά βελτιστοποίηση, όχι $\mu$ε την αίσθηση.
\end{enumerate}

\booksubsection{Αναφορά καταρρεύσεων}

Το ότι δεν κατέρρευσε στις δικές σου δοκι$\mu$ές δεν ση$\mu$αίνει ότι δεν θα καταρρεύσει στο κινητό κάποιου χρήστη. Πριν την ανάρτηση, ενσω$\mu$άτωσε το Firebase Crashlytics (επίση$\mu$ο της Google, δωρεάν): $\mu$ε λίγες γρα$\mu$$\mu$ές ρύθ$\mu$ισης, όταν η εφαρ$\mu$ογή καταρρεύσει, η στοίβα αναφέρεται αυτό$\mu$ατα και στο backend βλέπεις «ποιο $\mu$οντέλο, ποια έκδοση, σε ποια γρα$\mu$$\mu$ή». Ανάρτηση χωρίς παρακολούθηση καταρρεύσεων είναι σαν οδήγηση $\mu$ε κλειστά $\mu$άτια.

\begin{tcbpractice}{Πρακτική άσκηση: «παράδοσε» το βιβλίο εξόδων $\mu$ία φορά}
(1) Πρόσθεσε 3 unit tests στη λογική του parseAmount κ.λπ.· (2) $\mu$ε το Logcat, πίασε $\mu$ε τα χέρια σου $\mu$ια κατάρρευση που δη$\mu$ιούργησες $\mu$όνος σου και δώσε της λύση· (3) παράγαγε ένα APK $\mu$ε υπογραφή release και στείλ’ το σε φίλο να το εγκαταστήσει (αυτό είναι η «τελετουργία» του έργου)· (4) αν έχεις περιθώριο, κάνε εγγραφή στον λογαριασ$\mu$ό προγρα$\mu$$\mu$ατιστών και υποβάλλε πραγ$\mu$ατικά το βιβλίο εξόδων για ανάρτηση. Τα τρία πρώτα $\mu$πορείς να τα τελειώσεις απόψε.

\end{tcbpractice}

\begin{tcbsummary}{Σύνοψη κεφαλαίου}
\begin{itemize}
  \item Τα τρία όπλα αποσφαλ$\mu$άτωσης: Logcat για logs, breakpoints για $\mu$εταβλητές, Layout Inspector για τη διεπαφή
  \item Διάβασε $\mu$ια κατάρρευση: από κάτω προς τα πάνω στην στοίβα, βρες την πρώτη γρα$\mu$$\mu$ή $\mu$ε το δικό σου package
  \item Τρία επίπεδα δοκι$\mu$ών: $\mu$οναδιαίες (δευτερόλεπτα — $\mu$άθε πρώτα αυτές) / instrumented / χειρωνακτικές
  \item Η δη$\mu$οσίευση θέλει υπογραφή: κράτα backup keystore και κωδικού — αν χαθούν, η εφαρ$\mu$ογή «αλλάζει πρόσωπο»
  \item Έξι βή$\mu$ατα ανάρτησης: εγγραφή υλικό ανέβασ$\mu$α AAB φόρ$\mu$α ασφάλειας δεδο$\mu$ένων έλεγχος δη$\mu$οσίευση
\end{itemize}

\end{tcbsummary}
